% from ega1-5-fr.tex
\oldpage[I]{128}
La question étant locale, on est ramené à démontrer la dernière assertion. On
sait que $\widetilde{\mathfrak{N}}$ est un $\mathscr{O}_X$-Module
quasi-cohérent (\hyperref[I.1.4.1-fr]{1.4.1}) et que sa fibre au point
$x\in X$ est l'idéal $\mathfrak{N}_x$ de l'anneau de fractions $B_x$~; tout
revient à prouver que le nilradical de $B_x$ est contenu dans
$\mathfrak{N}_x$, l'inclusion opposée étant évidente. Or, soit $z/s$ un
élément du nilradical de $B_x$, avec $z\in B$, $s\notin\mathfrak{j}_x$~; par
hypothèse, il existe un entier $k$ tel que $(z/s)^k=0$, ce qui signifie qu'il
existe $t\notin\mathfrak{j}_x$ tel que $tz^k=0$. On en conclut $(tz)^k=0$, et
par suite $z/s=(tz)/(ts)$ appartient à $\mathfrak{N}_x$.

Nous dirons que le $\mathscr{O}_X$-Module quasi-cohérent $\mathscr{N}$ ainsi
défini est le \emph{Nilradical} de la $\mathscr{O}_X$-Algèbre $\mathscr{B}$~;
nous noterons en particulier $\mathscr{N}_X$ le nilradical de
$\mathscr{O}_X$.

\begin{corollary}[5.1.2]
\label{I.5.1.2-fr}
Soit $X$ un préschéma~; le sous-préschéma fermé de $X$ défini par le faisceau
d'idéaux $\mathscr{N}_X$ est le seul sous-préschéma réduit
(\hyperref[0.4.1.4-fr]{0, 4.1.4}) de $X$ ayant $X$ pour espace sous-jacent~;
c'est aussi le plus petit sous-préschéma de $X$ ayant $X$ pour espace
sous-jacent.
\end{corollary}

Comme le faisceau structural du sous-préschéma fermé défini $Y$ par
$\mathscr{N}_X$ est $\mathscr{O}_X/\mathscr{N}_X$, il est immédiat que $Y$
est réduit et a $X$ pour espace sous-jacent, puisque
$\mathscr{N}_x\ne\mathscr{O}_x$ pour tout $x\in X$. Pour démontrer les autres
assertions, remarquons qu'un sous-préschéma $Z$ de $X$ ayant $X$ pour espace
sous-jacent est défini par un faisceau d'idéaux $\mathscr{J}$
(\hyperref[I.4.1.3-fr]{4.1.3}) tel que
$\mathscr{J}_x\ne\mathscr{O}_x$ pour tout $x\in X$. On peut se borner au cas
où $X$ est affine, soit $X=\operatorname{Spec}(A)$ et
$\mathscr{J}=\widetilde{\mathfrak{J}}$, où $\mathfrak{J}$ est un idéal de
$A$~; alors, pour tout $x\in X$, on a
$\mathfrak{J}\subset\mathfrak{j}_x$, donc $\mathfrak{J}$ est contenu dans
tous les idéaux premiers de $A$, c'est-à-dire dans leur intersection
$\mathfrak{N}$, nilradical de $A$. Cela prouve que $Y$ est le plus petit
sous-préschéma de $X$ ayant $X$ comme espace sous-jacent
(\hyperref[I.4.1.9-fr]{4.1.9})~; en outre, si $Z$ est distinct de $Y$, on a
nécessairement $\mathscr{J}_x\ne\mathscr{N}_x$ pour un $x\in X$ au moins, et
par suite (\hyperref[I.5.1.1-fr]{5.1.1}) $Z$ n'est pas réduit.

\begin{definition}[5.1.3]
\label{I.5.1.3-fr}
On appelle préschéma réduit associé à un préschéma $X$ et on note
$X_{\mathrm{red}}$ l'unique sous-préschéma réduit de $X$ ayant $X$ pour
espace sous-jacent.
\end{definition}

Dire qu'un préschéma $X$ est réduit signifie donc que
$X=X_{\mathrm{red}}$.

\begin{proposition}[5.1.4]
\label{I.5.1.4-fr}
Pour que le spectre premier d'un anneau $A$ soit un préschéma réduit (resp.
intègre) (\hyperref[I.2.1.7-fr]{2.1.7}), il faut et il suffit que $A$ soit un
anneau réduit (resp. intègre).
\end{proposition}

En effet, il résulte aussitôt de (\hyperref[I.5.1.1-fr]{5.1.1}) que la
condition $\mathscr{N}=(0)$ est nécessaire et suffisante pour que
$X=\operatorname{Spec}(A)$ soit réduit~; l'assertion relative aux anneaux
intègres est alors une conséquence de
(\hyperref[I.1.1.13-fr]{1.1.13}).

Comme tout anneau de fractions $\ne\{0\}$ d'un anneau intègre est intègre, il
résulte de (\hyperref[I.5.1.4-fr]{5.1.4}) que pour tout préschéma
\emph{localement intègre} $X$, $\mathscr{O}_x$ est un anneau \emph{intègre}
pour tout $x\in X$. La réciproque est vraie lorsque l'espace sous-jacent à
$X$ est \emph{localement noethérien}~: en effet $X$ est alors réduit, et si
$U$ est un ouvert affine de $X$, qui soit un espace noethérien, $U$ n'a qu'un
nombre fini de composantes irréductibles, donc son anneau $A$ n'a qu'un
nombre fini d'idéaux premiers minimaux
(\hyperref[I.1.1.14-fr]{1.1.14}). Si deux de ces composantes $U_i$ avaient un
point commun $x$, $\mathscr{O}_x$ aurait au moins deux idéaux premiers
minimaux distincts, donc ne serait pas intègre~; les $U_i$ sont par suite des
ouverts deux à deux disjoints, et chacun d'eux est donc intègre.

\begin{env}[5.1.5]
\label{I.5.1.5-fr}
Soit $f=(\psi,\theta):X\to Y$ un morphisme de préschémas~; l'homomor-
